\section{Examples and extensions}\label{sec:examples}

\subsection{Two moduli spaces on a double plane}
Let $\phi:X\to\PP^2$ be a double cover branched over a smooth sextic, with
$\operatorname{Pic}(X)=\ZZ H$ and $H=\phi^*\OO_{\PP^2}(1)$.
Then $H^2=2$ and $\Vol(X)=1$ in our normalization. Put
\[
 F_0=\phi^*T_{\PP^2}(-1).
\]
The Euler sequence gives $\rk F_0=2$, $c_1(F_0)=H$, $c_2(F_0)=2$ and
$v(F_0)=(2,H,1)$, of square $-2$.

We verify the stability needed in the examples. A destabilizing saturated
rank-one subsheaf would extend to a nonzero map $\OO_X(aH)\to F_0$
with $a\ge1$. Since $\phi_*\OO_X=\OO\oplus\OO(-3)$,
\[
 H^0(X,F_0(-aH))=
 H^0(\PP^2,T_{\PP^2}(-1-a))\oplus
 H^0(\PP^2,T_{\PP^2}(-4-a))=0
\]
by the Euler sequence. Thus $F_0$ is stable.
It is spherical: simplicity and Serre duality give
$\operatorname{ext}^1(F_0,F_0)=v(F_0)^2+2=0$.
It is also the unique stable sheaf of its Mukai vector. A second,
nonisomorphic stable sheaf of that vector would have no homomorphisms in
either direction, whereas Riemann--Roch would give Euler characteristic
two, a contradiction.

\begin{example}\label{ex:four}
Take $v=(2,H,0)$. Then $N=2$, $c_2=3$ and $\kappa=5/2$.
Every semistable sheaf is slope stable, since equality of the rank-one
slope with $H^2/2$ would require an integer multiple of $H$ with
intersection $H^2/2$.
If a sheaf has defect length $\ell$, its stable double dual has square
$2-4\ell\ge-2$, so $\ell\le1$.
The non-locally-free locus is the projective bundle of one-dimensional
quotients of $F_0$ over $X$, of dimension three. Kernels of these quotients
exist, and their moduli space is smooth of dimension four. Its bundle
locus is therefore nonempty. The Uhlenbeck morphism contracts each
$\PP^1$ fibre and adds precisely one copy of $X$ to $\M$.

The completion has volume $128\pi^8$. Its boundary is totally geodesic
with metric $4\pi^2g_X$, and at every boundary point the tangent cone is
$\RR^4\times(\RR^4/\{\pm1\})$, of density $1/2$.
This Gieseker moduli space is generally not fine: the algebraic Mukai
pairings have common divisor two. The conclusions above do not require
a universal sheaf on the coarse space.
\end{example}

\begin{example}\label{ex:eight}
Take $v=(2,H,-1)$. Then $N=4$, $c_2=4$ and $\kappa=7/2$.
Stable kernels of quotients $F_0\to T$ of length two give nonempty
Gieseker moduli. The stable residual bundles have either vector $(2,H,0)$ and one bubble,
or vector $(2,H,1)$ and two bubbles. The latter residual bundle is $F_0$.
The length-one non-locally-free locus in the Gieseker space has dimension
$4+2+1=7$, including its $\PP^1$ Quot fibre. The length-two locus has
dimension six. Both are smaller than eight, so the bundle locus is
nonempty. The Uhlenbeck boundary consists of
\[
 \M_H(2,H,0)\times X
 \quad\text{and}\quad \{F_0\}\times\Sym^2X.
\]
The completion has volume $(16\pi^4)^4/24$.
The open one-bubble stratum has volume $(16\pi^4)^3/2$.
At $(F_0,x+y)$ with $x\ne y$, the tangent cone is
$\RR^8\times(\RR^4/\{\pm1\})^2$ and its density is $1/4$.
At $(F_0,2x)$, its affine cone with Reeb action is
$\frM_0(2,2)$, of complex dimension eight and density $11/108$.
There is no further $\CC^2$ factor: the centre-of-mass coordinates already
belong to $\frM_0(2,2)$.

Here $\chi(E)=r+s=1$, so the moduli space is fine. Its nontrivial
Uhlenbeck contraction and Corollary~\ref{cor:nonextension} give the fine
example relevant to the extension assertion discussed in
Section~\ref{sec:canonical}.
\end{example}

\subsection{A volume bound without a rational K\"ahler direction}
The potential argument used an integral polarization ray. Its exact volume
formula nevertheless gives an upper bound for all hyperk\"ahler metrics
on the differentiable K3 surface.

\begin{proposition}\label{prop:all-metrics}
Let $r\ge2$, $\gcd(r,k)=1$, and $N=rk-r^2+1\ge2$.
For any hyperk\"ahler metric $g$ on a K3 surface, the moduli space of
irreducible $SU(r)$ anti-self-dual connections of charge $k$ satisfies
\begin{equation}\label{eq:all-metrics}
 \Vol(\M_k(g))\le\frac{(16\pi^4\Vol(X,g))^N}{N!}.
\end{equation}
The volume of the empty moduli space is zero. If the positive three-plane
of $g$ contains a rational K\"ahler direction whose induced complex
surface has Picard number one, equality holds whenever the moduli space
is nonempty.
\end{proposition}
\begin{proof}
For a rational K\"ahler direction and Picard number one, the Mukai vector
$(r,0,r-k)$ is primitive. The polarization is Gieseker generic: a
same-reduced-Hilbert-polynomial factor would have rank divisible by $r$
because $\gcd(r,k)=1$. The Kobayashi--Hitchin correspondence identifies
the irreducible ASD locus with the slope-stable locally free locus.
When nonempty, Theorem~\ref{thm:canonical} and rescaling give equality.

Such positive three-planes are dense among positive three-planes in the
K3 period domain: choose a rational positive vector close to one vector
of the plane and choose the orthogonal positive two-plane generically to
avoid every other integral hyperplane. The K3 period description gives
nearby hyperk\"ahler metrics of the same limiting volume.
For any irreducible ASD connection, the hyperk\"ahler Weitzenb\"ock formula
identifies its obstruction space with the tensor product of its parallel
trace-free endomorphisms and the three-dimensional space of parallel
self-dual forms. Irreducibility makes this space zero.
The parameter-dependent implicit-function theorem therefore embeds each
compact subset of $\M_k(g)$ into the moduli spaces for nearby metrics,
with smooth convergence of their $L^2$ tensors. Its volume is bounded by
the total volumes of those nearby spaces, each at most the right side of
\eqref{eq:all-metrics}. Exhaustion proves the inequality.
\end{proof}

This is a lower-semicontinuity argument for total moduli volume, not a
completion theorem for an irrational positive three-plane. It does not
assert that nonempty equality cases are dense. For instance, if
$r\le k<2r$, then $\chi(E)=2r-k>0$; a slope-stable bundle of slope zero
has neither sections nor dual sections, contradicting Riemann--Roch.
The corresponding irreducible ASD moduli spaces are empty.

\subsection{A conditional expansion on a resolution}
The following elementary calculation records precisely what is needed
for an exceptional-area expansion. Suppose $r=2$, $D$ is an exceptional
divisor, and
\[
 q(\Lambda,D)=0,\qquad q(D)=-2(N-1),\qquad D\cdot R=-2
\]
for a contracted rational curve $R$.
Assume also that $A_t=\Lambda-t[D]$ is K\"ahler for small $t>0$.
The area of $R$ is $a=2t$. Put $Q=q(\Lambda)/2=16\pi^4V$.
Fujiki's formula gives
\begin{equation}\label{eq:resolution-volume}
 \Vol(\Lambda)-\Vol(A_t)
 =\frac{Q^{N-1}}{4(N-2)!}\,a^2+O(a^4).
\end{equation}
Indeed $q(A_t)/2=Q-(N-1)t^2$ and
$\Vol(A_t)=(Q-(N-1)t^2)^N/N!$.
When \eqref{eq:s1-volume} applies, its coefficient is
$\Vol(S_1)/4$. The displayed Beauville--Bogomolov identities and the
K\"ahler-ray hypothesis are assumptions here. The fibre intersection
$D\cdot R=-2$ alone does not prove either the square of $D$ or the
K\"ahler property of $A_t$.

\subsection{Further questions}
The continuous potential of Theorem~\ref{thm:continuous-potential} applies
to polarized surfaces beyond K3. To obtain a global class and a Quillen
metric there, one must separately control normalization, boundary
codimension, determinant-line descent and the flat character.
For a stratum inclusion, smoothness of a sheaf moduli space alone does
not show that every elementary modification is a limit of bundles;
this density condition must be established in the chosen component.

For K3 surfaces, several further metric questions remain.
At collisions of bubble points, does the tangent metric equal the flat
ADHM quotient metric without the extra uniqueness assertion
\textup{(Q)}? Do deeper strata with stable or reducible residual bundles
have a totally geodesic metric structure? What rates accompany the
Gromov--Hausdorff convergence near nonorbifold singularities?
The explicit potential is continuous, but the present argument gives
neither its optimal modulus of continuity nor a curvature-gluing expansion.
Extending the exact volume and completion statements to all irrational
K\"ahler directions, and treating nonprimitive vectors without a smooth
symplectic resolution, requires further arguments rather than a change
of notation in the theorems above.
